% ============================================================
%  I. INTRODUCTION
% ============================================================
\section{Introduction}
\label{sec:introduction}

\IEEEPARstart{E}{mergency} assessment protocols---including the \ac{ABCDE}
survey~\cite{thim2012abcde} and Focused Assessment with Sonography in Trauma (\ac{FAST})
ultrasound~\cite{rozycki1996}---enable rapid triage of life-threatening
conditions at the point of care. In mass-casualty
incidents and disaster sites, however, trained personnel are often unavailable,
creating a critical bottleneck.

Legged platforms such as the Boston Dynamics Spot can traverse stairs, rubble,
and confined spaces inaccessible to wheeled robots, and can therefore reach
patients wherever they lie. Navigation alone, however, is insufficient:
emergency assessment requires physical interaction with the patient---checking
airway patency and respiratory effort, inspecting the skin for signs of poor
circulation, scanning the body surface for wounds, burns, and haemorrhages, and
performing ultrasound at specific anatomical targets. A manipulator arm extends the robot's reachable
workspace, enables multiple viewpoints of the same body region, and provides the
dexterity needed for compliant probe contact.

% ── Related Work ─────────────────────────────────────────────

Robot-assisted ultrasound has advanced over two decades on fixed-base
platforms. Comprehensive surveys~\cite{vonhaxthausen2021,li2021overview}
document the evolution from teleoperated systems to autonomous scanning, and
Hidalgo et al.~\cite{hidalgo2023} review current clinical applications. Key
technical advances include active-sensing end-effectors~\cite{ma2022} and
autonomous protocol execution~\cite{suligoj2021,farsoni2025autonomous}. 

\draft{The Exposure phase of \ac{ABCDE} requires systematic visual
inspection of the body surface for wounds, burns, and haemorrhages.
Open-vocabulary detectors such as Grounding
DINO~\cite{liu2023groundingdino} enable zero-shot injury detection without
task-specific training data. Combined with depth back-projection from the
RGB-D stream, these models produce metric 3D injury locations.}
\draft{To plan the exposure scan grid and the \ac{FAST} the robot must first estimate
the patient's 3D body pose.} Neural Localizer Fields
(NLF)~\cite{isarandi2024nlf} \draft{regress 24-joint \ac{SMPL}
skeletons from monocular RGB, providing full anatomical coverage beyond the
standard 17-keypoint \ac{COCO} format.}
%Addressing uncertainty in human-centered robotics,
%Casado et al.~\cite{casado2025navigating} proposed diffusion-based intention
%estimation for shared wheelchair control---a theme that similarly motivates
%\ac{TERESA}'s quality-driven perception loop.
Despite
this progress, all deployed systems share a fundamental limitation: the patient
must be pre-positioned within the arm's reachable workspace, confining
autonomous ultrasound and scanning to controlled clinical settings. A possible solution is achieved by the 
modern \ac{WBC}. This builds on operational space formulation~\cite{khatib1987},
manipulability analysis~\cite{yoshikawa1985}, and optimisation-based
control~\cite{wensing2023}. These principles have been extended to legged
manipulators through hierarchical \ac{MPC}~\cite{sleiman2021},
\ac{QP}-based formulations~\cite{li2022wbc}, holistic~\cite{haviland2022} and
Spot-specific grasping and manipulability enhancement~\cite{zimmermann2021,xie2024manip}.
Across this body of work, however, the manipulation target is assumed known or
static; no formulation accounts for the coupling between body posture and the
quality of perceptual data used to localise that target in the first place.


Active perception research has established that deliberate sensor motion
systematically outperforms passive observation~\cite{zeng2020,bai2021,zaky2020}.
This insight has been operationalised through perception-aware
planning~\cite{zhang2020fisher,tordesillas2022}, next-best-view
grasping~\cite{breyer2022}, view planning for pose estimation~\cite{hu2022}, and
active neural sensing~\cite{ren2023}. When applied to mobile manipulation,
whole-body interaction~\cite{mittal2022}, coupled
perception-manipulation~\cite{xie2024capm} and one-shot multimodal
frameworks~\cite{qin2026} have demonstrated trajectory optimisation for
information gain. Critically, all of these systems operate on wheeled bases
whose body posture is fixed: the base provides mobility but does not participate
in shaping the perceptual field. Conversely, work on posture optimisation for
legged systems---through reachability
analysis~\cite{makhal2018,zhao2025}, base placement~\cite{fan2021},
body-ground contact~\cite{wolfslag2020}, and the broader legged manipulation
literature~\cite{gong2023}---treats perception as a static input to the motion
planner rather than as a quantity to be optimised. No existing system closes the
perception-action loop at the whole-body level on a legged platform for medical
tasks.

% ── Gap & Challenge ──────────────────────────────────────────

We advocate a paradigm in which body posture on a legged
platform is treated as a first-class perceptual resource: the robot actively
adapts its height, pitch, and yaw to maximise the quality of its own
perceptual input, rather than solely to minimise kinematic error. The key
idea is to close the perception-action loop at the whole-body level through
confidence-gated search and anticipatory scanning.

% ── TERESA ───────────────────────────────────────────────────

This paper presents \acf{TERESA} (Fig.~\ref{fig:system_overview}), an integrated
system for autonomous emergency assessment on the Boston Dynamics Spot legged
robot equipped with a Unitree Z1 manipulator arm --- \draft{Spotzi}.
Building on the fixed-base ultrasound scanning competence
established in~\cite{farsoni2025autonomous}, \ac{TERESA} extends autonomous
assessment to a fully mobile platform. \draft{We target the Exposure
and \ac{FAST} phases of \ac{ABCDE} because they share the same hardware
configuration and impose nearly identical whole-body coordination
demands.} \ac{TERESA} addresses this through a unified whole-body active
perception framework that treats body posture as an explicit control resource
across all four phases.

% ── Contributions ────────────────────────────────────────────

The key contributions of this paper are as follows.
\begin{enumerate}
  \item \textbf{Whole-body active perception framework} in which body
        posture is optimised to improve perceptual quality rather than
        solely to minimise kinematic error
        (Section~\ref{sec:active_perception}).
  \item \textbf{Confidence-gated search with dual-sensor locking} that
        combines forward base search with body pitch sweep, ensuring robust
        patient localisation before committing to navigation.
  \item \textbf{Exposure body scanning pipeline} integrating an adaptive
        skeleton-keypoint grid with zero-shot Grounding DINO injury
        detection and an interactive web-based review interface
        (Sections~\ref{sec:active_perception}
        and~\ref{sec:system_architecture}).
  \item \textbf{Pre-planned body reconfiguration} that searches over body
        height, pitch, and base displacement to place each anatomical target
        in the arm's preferred manipulation region
        (Section~\ref{sec:body_reconf}).
  \item \textbf{Experimental validation} on the \draft{Spotzi}
        platform demonstrating fully autonomous execution of all four phases
        (Section~\ref{sec:experimental_results}).
\end{enumerate}

To the best of our knowledge, \ac{TERESA} is the first legged robotic system to
demonstrate perception-driven whole-body reconfiguration for autonomous
emergency assessment.

% ── Outline ──────────────────────────────────────────────────

\draft{The remainder is organised as follows:
Section~\ref{sec:problem_formulation} formulates the problem,
Section~\ref{sec:active_perception} presents the whole-body active perception
framework, Section~\ref{sec:system_architecture} details the system
architecture and mission \ac{FSM},
Section~\ref{sec:experimental_results} reports experimental results, and
Section~\ref{sec:conclusion} concludes.}
